\section{Erd\H{o}s Problem \#43 (Round 5)}

\subsection{1) ROUND-5 OBJECTIVE}

\textbf{Path chosen: (C) --- obstruction/correction, with full proof of the corrected asymptotic statement.}

Round~4 already produced an explicit infinite-family disproof of the \emph{second} (equal-size) question, but only on an infinite subsequence and without extracting the full asymptotic consequence for \emph{all} $N$.  The most promising Round~5 direction is therefore to turn that disproof into a precise corrected theorem.

The new goal is to prove that the correct equal-size asymptotic is
\[
M(N):=\max\bigl\{m:\exists\text{ Sidon }A,B\subset[N],\ |A|=|B|=m,\ (A-A)\cap(B-B)=\{0\}\bigr\}
=\left(\frac{1}{\sqrt2}+o(1)\right)\sqrt N,
\]
and then to extend this to the unrestricted extremal quantity
\[
g(N):=\max\left\{\binom{|A|}{2}+\binom{|B|}{2}: A,B\subset[N]\ \text{Sidon},\ (A-A)\cap(B-B)=\{0\}\right\}.
\]
This will show that the second question is false in the strongest possible asymptotic sense, and that the first question is asymptotically correct up to an $o(N)$ error, leaving only the fine $O(1)$ scale unresolved.

\subsection{2) Round-4 FOUNDATION USED}

I explicitly use the following Round~4 results.
\begin{itemize}
\item \textbf{Round~4 parity-splitting construction.} For every odd prime power $q$, with
\[
N_q:=\frac{q^2-1}{2},
\]
there exist Sidon sets $A_q,B_q\subset[N_q]$ with $|A_q|=|B_q|=:m_q$, $(A_q-A_q)\cap(B_q-B_q)=\{0\}$, and
\[
m_q\ge \frac{q-\sqrt{2q-3}}{2} = \frac q2-O(\sqrt q).
\]
\item \textbf{Round~3 corrected convolution inequality.} If $A,B\subset[N]$ are Sidon and $(A-A)\cap(B-B)=\{0\}$, then for every integer $H\ge 1$,
\[
\frac{(|A|^2+|B|^2)H^2}{N+H}\le (|A|+|B|)H+H^2.
\]
\item \textbf{Round~4 status correction.} The second question is false; what remained open after Round~4 was to identify the \emph{correct} asymptotic statement and to advance the first question beyond finite computations.
\end{itemize}

\subsection{3) NEW INSIGHT / TOOL (ROUND-5)}

There are three genuinely new ingredients in this round.
\begin{itemize}
\item A \textbf{strengthened upper bound} that extends the Round~4 equal-size estimate to \emph{arbitrary} Sidon pairs:
\[
g(N)\le \frac N2+O(N^{3/4}).
\]
\item A \textbf{monotonicity + prime-number-theorem upgrade} of the Round~4 prime-power construction, which promotes the lower bound from an infinite subsequence to \emph{all sufficiently large $N$}.
\item Combining these yields a full corrected asymptotic theorem:
\[
M(N)=\left(\frac1{\sqrt2}+o(1)\right)\sqrt N,
\qquad
D_{\mathrm{eq}}(N):=2\binom{M(N)}2=\left(\frac12+o(1)\right)N,
\]
and also
\[
g(N)=\left(\frac12+o(1)\right)N=\binom{f(N)}2+o(N).
\]
\end{itemize}

\subsection{4) ATTACK PLAN (ROUND-5)}

\textbf{Exact gap left after Round~4.}
Round~4 settled the equal-size question only on an infinite subsequence and did not extract a theorem for the unrestricted extremal function $g(N)$.  The remaining tasks are:
\begin{enumerate}
\item prove a nontrivial upper bound for \emph{all} feasible pairs $(A,B)$, not just the equal-size case;
\item upgrade the Round~4 lower construction from the prime-power subsequence $N=(q^2-1)/2$ to all large $N$;
\item deduce the exact asymptotic scale of both $M(N)$ and $g(N)$.
\end{enumerate}

\textbf{Why this overcomes the earlier obstacle.}
The Round~3 corrected inequality is already strong enough to force an $N/2+O(N^{3/4})$ upper bound once one controls $|A|+|B|$.  Meanwhile the Round~4 parity-splitting construction already gives near-extremizers on a dense-enough subsequence; the prime number theorem allows us to pull those constructions down to arbitrary large $N$ by monotonicity.

\subsection{5) WORK (ROUND-5)}

\subsubsection{5.1 Definitions and monotonicity}

Define
\[
M(N):=\max\bigl\{m:\exists\text{ Sidon }A,B\subset[N],\ |A|=|B|=m,\ (A-A)\cap(B-B)=\{0\}\bigr\},
\]
and
\[
g(N):=\max\left\{\binom{|A|}{2}+\binom{|B|}{2}: A,B\subset[N]\ \text{Sidon},\ (A-A)\cap(B-B)=\{0\}\right\}.
\]
Both $M(N)$ and $g(N)$ are nondecreasing in $N$, because any feasible pair inside $[N']$ is also feasible inside $[N]$ whenever $N'\le N$.

For later use, note that if $S\subset[N]$ is Sidon and $|S|=s$, then its positive differences are distinct elements of $\{1,\dots,N-1\}$, so
\[
\binom{s}{2}\le N-1.
\]
Hence every Sidon subset of $[N]$ has size
\begin{equation}\label{eq:sidon-size-trivial}
s\le \frac{1+\sqrt{8N-7}}{2}=\sqrt{2N}+O(1).
\end{equation}

\subsubsection{5.2 A new upper bound for arbitrary pairs}

\begin{proposition}\label{prop:g-upper}
Let $A,B\subset[N]$ be Sidon sets with $(A-A)\cap(B-B)=\{0\}$.  Write
\[
a:=|A|,\qquad b:=|B|,\qquad S:=a^2+b^2,\qquad T:=a+b.
\]
Then
\[
S\le N+O(N^{3/4}).
\]
Consequently,
\[
\binom{|A|}{2}+\binom{|B|}{2}\le \frac N2+O(N^{3/4}),
\]
and therefore
\[
g(N)\le \frac N2+O(N^{3/4}).
\]
\end{proposition}

\begin{proof}
We use the Round~3 corrected convolution inequality:
\[
\frac{SH^2}{N+H}\le TH+H^2
\qquad (H\ge 1,\ H\in\mathbb Z).
\]
Rearranging gives
\begin{equation}\label{eq:S-basic}
S\le \frac{(T+H)(N+H)}{H}=N+T+H+\frac{TN}{H}.
\end{equation}

By \eqref{eq:sidon-size-trivial}, each of $a,b$ is $O(\sqrt N)$, hence
\begin{equation}\label{eq:T-upper}
T=O(\sqrt N).
\end{equation}
If $T=0$ there is nothing to prove.  Otherwise choose
\[
H:=\left\lceil\sqrt{TN}\right\rceil.
\]
Then
\[
H\le \sqrt{TN}+1,
\qquad
\frac{TN}{H}\le \sqrt{TN}.
\]
Substituting into \eqref{eq:S-basic},
\[
S\le N+T+2\sqrt{TN}+1.
\]
Using \eqref{eq:T-upper}, we have $T=O(\sqrt N)$ and therefore $\sqrt{TN}=O(N^{3/4})$.  Hence
\[
S\le N+O(N^{3/4}).
\]
Finally,
\[
\binom{|A|}{2}+\binom{|B|}{2}
=\frac{a(a-1)+b(b-1)}{2}
=\frac{S-T}{2}
\le \frac N2+O(N^{3/4}),
\]
proving the proposition.
\end{proof}

\begin{corollary}\label{cor:M-upper}
For the equal-size extremal function one has
\[
M(N)\le \frac1{\sqrt2}\sqrt N+O(N^{1/4}).
\]
\end{corollary}

\begin{proof}
If $|A|=|B|=m$, then Proposition~\ref{prop:g-upper} gives
\[
2m^2\le N+O(N^{3/4}),
\]
whence
\[
m\le \frac1{\sqrt2}\sqrt N+O(N^{1/4}).
\qedhere
\]
\end{proof}

\subsubsection{5.3 Upgrading the Round~4 lower bound to all large $N$}

The Round~4 construction gives large equal-size pairs only for $N_q=(q^2-1)/2$ with $q$ an odd prime power.  To pass to arbitrary $N$, we use monotonicity together with the prime number theorem.

\begin{lemma}\label{lem:prime-below}
For every real $x\to\infty$, there exists an odd prime $q\le x$ such that
\[
q=(1-o(1))x.
\]
\end{lemma}

\begin{proof}
Fix $\varepsilon\in(0,1)$.  By the prime number theorem,
\[
\pi(x)-\pi((1-\varepsilon)x)\sim \frac{\varepsilon x}{\log x}\to\infty
\qquad (x\to\infty).
\]
Therefore, for all sufficiently large $x$, the interval $[(1-\varepsilon)x,x]$ contains a prime.  Since $\varepsilon>0$ is arbitrary, there exists a prime $q\le x$ with $q=(1-o(1))x$.  For large $x$ this prime is automatically odd.
\end{proof}

\begin{proposition}\label{prop:M-lower-allN}
One has
\[
M(N)\ge \left(\frac1{\sqrt2}-o(1)\right)\sqrt N.
\]
In particular,
\[
M(N)=\left(\frac1{\sqrt2}+o(1)\right)\sqrt N.
\]
\end{proposition}

\begin{proof}
Set
\[
x:=\sqrt{2N+1}.
\]
By Lemma~\ref{lem:prime-below}, choose an odd prime $q\le x$ such that $q=(1-o(1))x$.  Since every prime is a prime power, we may apply the Round~4 parity-splitting construction at this $q$.  Let
\[
N_q:=\frac{q^2-1}{2}\le \frac{x^2-1}{2}=N.
\]
Round~4 then supplies Sidon sets $A_q,B_q\subset[N_q]$ with
\[
|A_q|=|B_q|=:m_q\ge \frac{q-\sqrt{2q-3}}{2} = \frac q2-O(\sqrt q).
\]
By monotonicity of $M$,
\[
M(N)\ge M(N_q)\ge m_q.
\]
Because $q=(1-o(1))x=(1-o(1))\sqrt{2N+1}$ and $\sqrt q=o(\sqrt N)$, we conclude that
\[
M(N)\ge \frac q2-o(\sqrt N)
=\left(\frac1{\sqrt2}-o(1)\right)\sqrt N.
\]
Combining this with Corollary~\ref{cor:M-upper} yields
\[
M(N)=\left(\frac1{\sqrt2}+o(1)\right)\sqrt N.
\]
\end{proof}

\subsubsection{5.4 The corrected asymptotic theorem}

\begin{theorem}[Equal-size case]
\label{thm:eq-corrected}
Let
\[
D_{\mathrm{eq}}(N):=2\binom{M(N)}{2}=M(N)(M(N)-1).
\]
Then
\[
D_{\mathrm{eq}}(N)=\left(\frac12+o(1)\right)N.
\]
Equivalently, the sharp asymptotic constant in the equal-size problem is $1/\sqrt2$, and there is no fixed positive gap below it.
\end{theorem}

\begin{proof}
By Proposition~\ref{prop:M-lower-allN},
\[
M(N)=\left(\frac1{\sqrt2}+o(1)\right)\sqrt N.
\]
Therefore
\[
D_{\mathrm{eq}}(N)=M(N)(M(N)-1)=\left(\frac12+o(1)\right)N.
\]
This proves the claim.
\end{proof}

\begin{theorem}[Unrestricted case]
\label{thm:g-corrected}
One has
\[
g(N)=\left(\frac12+o(1)\right)N.
\]
Since $f(N)\sim \sqrt N$, this is equivalent to
\[
g(N)=\binom{f(N)}{2}+o(N).
\]
\end{theorem}

\begin{proof}
The lower bound follows from the equal-size case:
\[
g(N)\ge D_{\mathrm{eq}}(N)=\left(\frac12+o(1)\right)N.
\]
The upper bound is Proposition~\ref{prop:g-upper}, namely
\[
g(N)\le \frac N2+O(N^{3/4})=\left(\frac12+o(1)\right)N.
\]
Hence
\[
g(N)=\left(\frac12+o(1)\right)N.
\]
Finally, the classical asymptotic $f(N)\sim \sqrt N$ implies
\[
\binom{f(N)}2=\left(\frac12+o(1)\right)N,
\]
so indeed
\[
g(N)=\binom{f(N)}2+o(N).
\qedhere
\]
\end{proof}

\subsection{6) ADVERSARIAL VERIFICATION}

\paragraph{(i) Did Round~5 improperly reprove Round~4 material?}
No.  The Round~4 parity-splitting construction is used as a black box in Proposition~\ref{prop:M-lower-allN}.  The only proof repeated from earlier rounds is the convolution-based upper estimate, and there it is \emph{strengthened} from the equal-size setting to arbitrary pairs, which is explicitly permitted.

\paragraph{(ii) Hidden dependence on prime powers?}
There is none.  Proposition~\ref{prop:M-lower-allN} uses primes, and every prime is a prime power, so the Round~4 construction applies unchanged.

\paragraph{(iii) Monotonicity check.}
If $A,B\subset[N_q]$ are feasible, then $A,B\subset[N]$ for every $N\ge N_q$.  Thus both $M(N)$ and $g(N)$ are indeed nondecreasing.  No rescaling or translation is required.

\paragraph{(iv) Prime-number-theorem quantifiers.}
Lemma~\ref{lem:prime-below} is uniform in the needed sense: for every fixed $\varepsilon>0$ and all sufficiently large $x$, a prime exists in $[(1-\varepsilon)x,x]$.  This is exactly what is needed to make the multiplicative loss $o(1)$.

\paragraph{(v) Does Proposition~\ref{prop:g-upper} really give an $o(N)$ error?}
Yes.  The proof yields $g(N)\le N/2+O(N^{3/4})$, and $N^{3/4}=o(N)$.  Hence Theorem~\ref{thm:g-corrected} is rigorous.

\paragraph{(vi) What remains open after Round~5?}
The original first question asked for the much sharper estimate
\[
g(N)\le \binom{f(N)}2+O(1).
\]
Round~5 proves only
\[
g(N)=\binom{f(N)}2+o(N),
\]
so the fine error term remains unresolved.  The second question, however, is fully settled negatively, and its corrected asymptotic form is now proved.

\subsection{7) FINAL}

\textbf{FULL SOLUTION --- FULL PROOF.}

The second question is false, and the corrected asymptotic statement is now proved in full:
\[
M(N)=\left(\frac1{\sqrt2}+o(1)\right)\sqrt N,
\qquad
D_{\mathrm{eq}}(N)=\left(\frac12+o(1)\right)N.
\]
Equivalently, in the equal-size case the sharp asymptotic constant is exactly $1/\sqrt2$, so no fixed positive gap $c>0$ can occur.

Moreover, for the unrestricted extremal quantity,
\[
g(N)=\left(\frac12+o(1)\right)N=\binom{f(N)}2+o(N).
\]
Thus the \emph{asymptotic} form of the first question is true, but the original $O(1)$ refinement remains open.

\subsection{8) COMPLETION ESTIMATE (MANDATORY)}

\textbf{COMPLETION: 90\%}

\subsection{9) REFERENCES}

\begin{itemize}
\item Round~4 write-up supplied in this task, especially the parity-splitting Bose--Chowla construction for $N=(q^2-1)/2$.
\item M. B. Nathanson, \emph{The Bose--Chowla argument for Sidon sets}, J. Number Theory 238 (2022), 133--146.
\item B. Green, \emph{The number of squares and $B_h[g]$ sets}, especially the summary of the Erd\H{o}s--Tur\'an upper bound and Singer lower bound implying $f(N)\sim \sqrt N$.
\item The prime number theorem (standard).
\item Erd\H{o}s Problems, Problem \#43 page and discussion thread (current status of the two questions).
\end{itemize}
